\sscp{{\#}kisor(o)}{14-03-03}
Forms similar to \tsc{North Halmahera} *kusor occur in half a
dozen languages belonging to different language families on the
Bird’s Head and Onin peninsula of Papua,
as well as the Raja Ampat Islands.
Examples are given in \trf{14-05}.
from \citet[263]{SmitsVoorhoeve1992b}[118, 219, 221, 228f], \citeyear{SmitsVoorhoeve1998}.

\begin{table}[h]\centering\tcp{Words resembling *kusor/{\#}kisor(o)}{14-05}
%\begin{adjustbox}{width=\textwidth}
	\begin{threeparttable}\st{2pt}
	\st{5.5pt}\begin{tabular}{l|llll}\lsptoprule
Family	&	Language	&	ISO	&	form	&	gloss\su{†}\\	\midrule
\tsc{Austronesian {(STB)}}	&	Onin	&	oni	&	kisor	&	tree wallaby\\
\tsc{West Bird's Head}	&	Mooi	&	mxn	&	kasyor	&	tree wallaby\\
\tsc{East Bird's Head}	&	Sougb	&	mnx	&	ucira, usira	&	marsupial\\
\tsc{Konda-Yahadian\su{‡} }	&	Yahadian	&	ner	&	kisoro	&	kangaroo\\
\tsc{Inanwatan-Duriankere\su{‡}}	&	Duriankere	&	dbn	&	kesyoro	&	tree wallaby\\
\tsc{Inanwatan-Duriankere\su{‡}}	&	Inanwatan	&	szp	&	isoro	&	tree wallaby\\
\tsc{South Bird's Head\su{‡}}	&	Kais	&	kzm	&	kisoro	&	tree wallaby\\
		\lspbottomrule
	\end{tabular}
		\begin{tablenotes}\tnsize
			\item[†] {Glosses for all forms in this table are not certain.}
			\item[‡] {\tsc{Konda-Yahadian} and \tsc{Inanwatan-Duriankere}
								may belong to a larger \tsc{South Bird's Head}
								group which has weak but disputed claims to belong
								to \tsc{Trans New Guinea} \citep[75]{PawleyHammarstrom2018}.}
		\end{tablenotes}
	\end{threeparttable}
%\end{adjustbox}
\end{table}

Further morphological,
phonological, or comparative analysis of the forms presented
in \trf{14-05} is not currently possible due to the lack of detailed
descriptions of most of the languages.
The similarity of the forms across three to five unrelated language
families points to these forms being spread by contact.
Of these languages, Inanwatan has been described by \citet{deVries2004},
and some comments on the form \it{isoro} ‘wallaby,
tree wallaby’ can be made on the basis of this description.

Firstly, Inanwatan has a contrast between words with an initial
glottal stop and an initial vowel.
The form \it{isoro} ‘(tree) wallaby’ may thus have an initial
contrastive glottal stop which is not transcribed in the lists
presented by \citet[219, 229]{SmitsVoorhoeve1998}.
Secondly, Inanwatan does not allow final consonants,
and “vowel insertion and consonant deletion are employed to adapt
foreign words to Inanwatan phonotactics” \citep[6]{deVries2004}.
Examples include Malay \it{wakil} > Inanwatan [ˈmakiri] ‘deputy
headman’; Dutch \it{emmer} > Malay \it{ember} > Inanwatan [ˈɛβɛro]
‘bucket’; and Malay \it{kapal} > [ˈkaparo] ‘ship’.
Thus, if \it{isoro} has been acquired by contact the final vowel
may be an insertion.
Finally, most nouns in Inanwatan end in final [o] \citep[33]{deVries2004}.
Thus, \it{isoro} ‘(tree) wallaby’ fits into the expected patterns of the language.